% Modelado de Intrusiones con Grafos Temporales — GraphHunter
% Reconstrucción LaTeX del paper JAIIO, actualizada al estado del motor
% de agosto de 2026. Compilar con: pdflatex (x2, por las referencias).
\documentclass[11pt,a4paper]{article}

\usepackage[utf8]{inputenc}
\usepackage[T1]{fontenc}
\usepackage[spanish,es-noshorthands]{babel}
\usepackage{amsmath,amssymb}
\usepackage{booktabs}
\usepackage{geometry}
\geometry{margin=2.6cm}
\usepackage{url}
\usepackage[hidelinks]{hyperref}
\usepackage{enumitem}

\newcommand{\code}[1]{\texttt{#1}}

\title{\vspace{-1.2cm}\textbf{Modelado de Intrusiones con Grafos Temporales:\\
Optimizando el Threat Hunting Proactivo}}
\author{Lucas Sotomayor\thanks{BASE4 Security, Buenos Aires, Argentina.
\code{lsotomayor@base4sec.com}} \and
Diego Staino\thanks{BASE4 Security, Buenos Aires, Argentina.
\code{dstaino@base4sec.com}}}
\date{}

\begin{document}
\maketitle

% ======================================================================
\subsection*{Resumen}
\small
Este trabajo propone un modelo de grafo de conocimiento temporal tipado para
telemetría de seguridad y un algoritmo de búsqueda en profundidad dirigido por
hipótesis para Threat Hunting proactivo. A diferencia de las plataformas SIEM
tradicionales, que almacenan la telemetría en formato tabular y requieren
pivoteo iterativo manual para reconstruir cadenas de ataque, el enfoque
propuesto materializa las relaciones causales entre eventos de seguridad como
un multigrafo dirigido con diez tipos de entidades y once tipos de relaciones.
Los analistas codifican hipótesis de ataque como cadenas tipadas secuenciales,
y el sistema retorna todos los caminos temporalmente válidos mediante reglas de
poda integradas (tipo de relación, tipo de entidad, monotonía causal con
tolerancia de desfasaje de relojes, ventana temporal, huecos acotados y
$k$-simplicidad). La evaluación empírica sobre datasets reales ---incluyendo
telemetría Sysmon, logs de Azure Sentinel y el corpus Mordor (963K
eventos)--- demuestra aceleraciones de $1{,}747\times$ a $11{,}600\times$
respecto al baseline sin poda. Un barrido de escala a densidad constante mide
un crecimiento de la latencia cercano a $O(|E|^{1{,}2})$, con una mediana de
$121$~ms sobre $2{,}5$~millones de aristas. La comparativa operativa muestra
una reducción de 8--20 minutos y múltiples consultas manuales por investigación
(SIEM) a 1--2 minutos con una sola formulación de hipótesis.

\paragraph{Palabras clave:} Threat Hunting, grafo de conocimiento temporal,
telemetría de seguridad, búsqueda dirigida por hipótesis, MITRE ATT\&CK.

% ======================================================================
\subsection*{\textit{Intrusion Modeling with Temporal Graphs: Optimizing
Proactive Threat Hunting}}

\textbf{Abstract.} This work proposes a temporal typed knowledge graph model
for security telemetry and a hypothesis-driven depth-first search algorithm for
proactive Threat Hunting. Unlike traditional SIEM platforms, which store
telemetry in tabular form and require iterative manual pivoting to reconstruct
attack chains, the proposed approach materializes causal relationships between
security events as a directed multigraph with ten entity types and eleven
relation types. Analysts encode attack hypotheses as typed sequential chains,
and the system returns all temporally valid paths through integrated pruning
rules (relation type, entity type, causal monotonicity with clock-skew
tolerance, time window, bounded gaps, and $k$-simplicity). Empirical evaluation
on real-world datasets ---including Sysmon telemetry, Azure Sentinel logs, and
the Mordor corpus (963K events)--- demonstrates speedups from $1{,}747\times$
to $11{,}600\times$ over unpruned baselines. A constant-density scaling sweep
measures latency growth close to $O(|E|^{1.2})$, with a $121$~ms median over
$2.5$ million edges. The operational comparison shows a reduction from 8--20
minutes and multiple manual queries per investigation (SIEM) to 1--2 minutes
with a single hypothesis formulation, representing an order-of-magnitude
decrease in Mean Time to Investigate (MTTI).

\paragraph{Keywords:} Threat Hunting, temporal knowledge graph, security
telemetry, hypothesis-driven search, MITRE ATT\&CK.
\normalsize

% ======================================================================
\section{Introducción}

\subsection{Threat Hunting: contexto y necesidad}
El Threat Hunting (caza de amenazas) es una disciplina proactiva que consiste
en la búsqueda iterativa, basada en hipótesis, de amenazas avanzadas que han
evadido controles automatizados (Hutchins et al., 2011). El proceso sigue un
ciclo: formulación de hipótesis sobre comportamiento adversarial, análisis de
registros, identificación de patrones anómalos y retroalimentación hacia
controles defensivos. Marcos como MITRE ATT\&CK\footnote{\url{https://attack.mitre.org}}
proporcionan una taxonomía estandarizada de tácticas, técnicas y procedimientos
(TTPs) que estructuran este proceso y permiten a los analistas formular
hipótesis con base en comportamientos adversariales conocidos.

La efectividad del Threat Hunting depende de dos factores críticos: la calidad
de los registros disponibles y la capacidad de las herramientas/personas para
revelar relaciones causales entre eventos. En este segundo aspecto, las
herramientas predominantes presentan limitaciones estructurales significativas.

\subsection{Planteo del problema: divergencia estructural}
Existe una diferencia fundamental entre cómo se almacena la telemetría de
seguridad y cómo operan los adversarios. Los SIEM (Security Information and
Event Management) almacenan registros en formato tabular optimizado para
indexación y búsqueda por campos, mientras que los adversarios ejecutan cadenas
causales de acciones: compromiso inicial, escalada de privilegios, movimiento
lateral y exfiltración.

Las plataformas modernas ofrecen dashboards, alertas estadísticas y capacidades
de análisis de comportamiento de usuarios (UEBA), pero estas agregaciones
operan a nivel macro y omiten la causalidad a nivel micro entre eventos
individuales. Esta divergencia genera tres barreras operativas concretas para
el analista de SOC:

\begin{enumerate}[leftmargin=*]
  \item \textbf{Pivoteo iterativo ineficiente.} Reconstruir una cadena de
  ataque de $L$ pasos requiere aproximadamente $L$ consultas manuales
  secuenciales (2--5 minutos cada una), con un tiempo total de 8--20 minutos
  para $L = 4$. Cada consulta requiere la extracción manual de identificadores
  (IPs, nombres de usuario, PIDs) del resultado anterior para formular la
  siguiente.
  \item \textbf{Sobrecarga cognitiva.} Al carecer de una representación visual
  de relaciones, el analista debe mantener mentalmente un ``grafo'' de
  dependencias causales entre eventos mientras navega resultados. La capacidad
  de la memoria de trabajo humana ($7 \pm 2$ ítems según Miller, 1956) se
  convierte en un cuello de botella cuando las cadenas superan los 3--4 pasos.
  \item \textbf{Fragmentación del contexto.} El volumen de telemetría y el
  formato lineal de los logs provocan pérdida de contexto. Eventos causalmente
  relacionados quedan dispersos entre miles de registros irrelevantes,
  dificultando la visión integral del incidente.
\end{enumerate}

\subsection{Propuesta y objetivos}
Este trabajo propone la transición del análisis tabular a un modelo de grafo
temporal tipado que materializa las relaciones causales entre eventos de
seguridad sobre un sistema completo ---de la ingestión de logs a la búsqueda de
patrones--- donde el analista plantea hipótesis de caza como consultas de
``caminos tipados'' y obtiene resultados visuales. Los objetivos específicos
son: (1) definir un modelo de datos sobre un multigrafo dirigido, etiquetado y
temporal para telemetría heterogénea (Sysmon, Azure Sentinel, JSON, CSV,
HAProxy, DARPA TC CDM) alineada con MITRE ATT\&CK; (2) diseñar un algoritmo DFS
con podas que explote el tipado y la temporalidad del dominio para lograr
tiempo esperado lineal; (3) validar el método sobre datasets realistas
($> 900$K eventos) midiendo la reducción de MTTI frente al flujo SIEM;
(4) implementar una herramienta que integre ingestión, búsqueda, puntuación de
anomalías y visualización para SOCs; (5) integrar clasificación GNN sobre el
mismo grafo, cerrando el ciclo detección-investigación con aceleración opcional
por NPU/GPU vía DirectML; y (6) exponer las operaciones del motor mediante un
servidor MCP para asistentes de IA.

% ======================================================================
\section{Estado del arte}

\subsection{Herramientas SIEM y sus limitaciones para Threat Hunting}
La piedra angular de los Centros de Operaciones de Seguridad (SOC) han sido los
SIEM ---Splunk, Microsoft Sentinel, Elastic Security--- que almacenan
telemetría en formato tabular optimizado para búsqueda full-text y agregaciones
estadísticas. Si bien estas plataformas ofrecen capacidades de búsqueda,
dashboards y alertas, su diseño no materializa relaciones causales entre
eventos.

Para mitigar la fatiga de alertas, la industria incorporó análisis estadísticos
y heurísticos. Las soluciones UEBA (User and Entity Behavior Analytics)
detectan anomalías en patrones de comportamiento, proporcionando conciencia
situacional a nivel macro. Sin embargo, no resuelven el problema de la
reconstrucción causal micro: determinar qué proceso específico generó qué
conexión de red hacia qué servidor, y qué sucedió después.

Las capacidades de grafos en plataformas comerciales ---como los Investigation
Graphs de Microsoft Sentinel o los Causality Views de Palo Alto Cortex XDR---
representan un avance, pero operan como capas de visualización sobre datos
tabulares subyacentes. No explotan formalmente la estructura del grafo para
búsqueda de patrones, y las travesías con restricciones temporales típicamente
superan los 100~ms para más de 3 saltos.

\subsection{Enfoques de grafos en ciberseguridad}
BloodHound\footnote{\url{https://github.com/SpecterOps/BloodHound}} (Robbins et
al., 2017) demostró el poder de la modelización relacional en Active Directory,
permitiendo descubrir rutas de escalada de privilegios mediante travesía de
grafos de permisos. Sin embargo, opera sobre relaciones estáticas ---permisos y
membresías de grupos--- sin componente temporal, lo que lo limita a análisis de
configuración y no a detección de actividad maliciosa en tiempo real.

SLEUTH (Hossain et al., 2017) y HOLMES (Milajerdi et al., 2019a) construyen
grafos de proveniencia a partir de llamadas de sistema y aplican reglas de
detección basadas en flujos de información. Poirot (Milajerdi et al., 2019b)
alinea grafos de ataque provenientes de inteligencia de amenazas contra logs de
auditoría mediante alineación aproximada nodo-a-nodo. ProvDetector (Wang et
al., 2020) utiliza embeddings de caminos para detección de anomalías en grafos
de proveniencia. Unicorn (Han et al., 2020) emplea histosketches para detección
de APT a nivel de grafo completo.

Estos sistemas comparten tres rasgos que los distinguen de nuestro enfoque:
operan sobre logs de llamadas de sistema (no telemetría heterogénea
multi-fuente), emplean detección automática por reglas o modelos aprendidos
(sin hipótesis ad-hoc del analista), y no ofrecen una interfaz de búsqueda
interactiva. El más cercano, Poirot, también alinea grafos de ataque, pero
contra un grafo de consulta fijo derivado de CTI, mientras que nuestro enfoque
permite matching exacto de caminos temporales con hipótesis formuladas
dinámicamente, priorizando interpretabilidad y control humano.

\subsection{Subgrafo temporal y joins worst-case-optimal}
En el ámbito de algoritmos de grafos, Ti\&To (Kim et al., 2022) intercala
filtrado topológico y temporal para isomorfismo de subgrafos temporal, análogo
a nuestra poda integrada tipo y temporal. Sin embargo, la propuesta de este
trabajo especializa el problema a consultas de caminos tipados secuenciales,
permitiendo una estructura DFS lineal sin la sobrecarga del isomorfismo
general. El método edge-driven cronológico (Cai et al., 2023) procesa aristas
en orden temporal, y SAL/SAM (Lai et al., 2022) ofrece ejecución distribuida
streaming.

Una segunda familia relevante son los \emph{joins worst-case-optimal} (WCOJ).
Leapfrog Triejoin (Veldhuizen, 2014) y Generic Join enumeran en tiempo óptimo
respecto de la cota AGM sin materializar productos intermedios, lo que resulta
decisivo en patrones cíclicos donde el join binario explota; EmptyHeaded
(Aberger et al., 2017) combina WCOJ con intersección SIMD. Para consultas
acíclicas, Yannakakis (1981) enumera en $O(|D| + |\mathrm{OUT}|)$. El sistema
descrito aquí incorpora ambos como operadores físicos especializados junto al
DFS, con la salvedad ---documentada en la Sección~\ref{sec:operadores}--- de
que su semántica temporal debe alinearse explícitamente con la del DFS, so pena
de responder una pregunta distinta de la formulada.

\subsection{Detección basada en aprendizaje sobre grafos}
Trabajos recientes aplican Graph Neural Networks (GNN) a grafos de proveniencia
para detección automática de APT. ThreaTrace (Wang et al., 2022) entrena
modelos GNN para clasificar nodos en grafos de proveniencia como benignos o
maliciosos. MAGIC (Jia et al., 2024) utiliza aprendizaje de representaciones
enmascarado sobre grafos para detección de APT a nivel de sistema. Sin embargo,
estos sistemas operan como detectores de caja negra sin capacidad de
investigación interactiva: generan alertas pero no permiten al analista
formular hipótesis ad-hoc ni explorar el contexto causal de la alerta. Nuestra
propuesta cierra esta brecha integrando la clasificación GNN como un componente
de la puntuación de anomalía, de modo que la detección automática alimenta la
priorización de caminos en la búsqueda dirigida por hipótesis, y viceversa.

% ======================================================================
\section{Modelo de grafo para telemetría de seguridad}

\subsection{Definición de grafo de conocimiento temporal de amenazas}
\textbf{Grafo de conocimiento temporal de amenazas.} Una tupla
\[
G = (V, E, \tau, \sigma_V, \sigma_E, \mu_V, \mu_E)
\]
donde $V$ es un conjunto finito de vértices (entidades de seguridad); $E$ es
una bolsa (multiconjunto) sobre $V \times V$ con identificadores únicos;
$\tau : E \to \mathbb{Z}$ asigna marcas temporales; $\sigma_V : V \to \Sigma_V$
y $\sigma_E : E \to \Sigma_E$ asignan tipos; y $\mu_V, \mu_E$ asignan
diccionarios de metadatos.

\textbf{Alfabetos de tipos de seguridad.}
\begin{align*}
\Sigma_V &= \{\mathsf{IP}, \mathsf{Host}, \mathsf{User}, \mathsf{Process},
\mathsf{File}, \mathsf{Domain}, \mathsf{Registry}, \mathsf{Email},
\mathsf{URL}, \mathsf{Session}\}\\
\Sigma_E &= \{\mathsf{Auth}, \mathsf{Connect}, \mathsf{Execute}, \mathsf{Read},
\mathsf{Write}, \mathsf{DNS}, \mathsf{Spawn}, \mathsf{Modify}, \mathsf{Alert},
\mathsf{Uses}, \mathsf{From}\}
\end{align*}
con $|\Sigma_V| = 10$ y $|\Sigma_E| = 11$.

La elección de $E$ como multiconjunto es fundamental para el dominio: un mismo
usuario puede autenticarse al mismo host repetidas veces, cada evento con marca
temporal y contexto distintos. Esto distingue el modelo de enfoques estáticos
como BloodHound, donde una arista ``tiene-permiso'' es binaria. Como se
detallará en la Sección~\ref{sec:operadores}, las aristas paralelas no son un
detalle de representación: un operador que colapse cada par de entidades a una
única marca temporal descarta coincidencias válidas.

Estos alfabetos de tipos fueron diseñados para cubrir las entidades y
relaciones más frecuentes en telemetría de seguridad según la taxonomía MITRE
ATT\&CK. Cada tipo de relación mapea naturalmente a una o más técnicas ATT\&CK:
$\mathsf{Auth}$ cubre técnicas de acceso con credenciales válidas (T1078),
$\mathsf{Execute}$ y $\mathsf{Spawn}$ cubren ejecución (TA0002),
$\mathsf{Connect}$ y $\mathsf{DNS}$ cubren comunicaciones C2 (TA0011), y
$\mathsf{Read}$/$\mathsf{Write}$ cubren acceso a datos y exfiltración (TA0009,
TA0010).

\subsection{La sesión como entidad de primera clase}
\label{sec:sesion}
El tipo $\mathsf{Session}$ y las relaciones $\mathsf{Uses}$ y $\mathsf{From}$
merecen una justificación específica, porque su ausencia en el modelo original
resultó ser una limitación material observada en despliegue.

Los identificadores de correlación (\code{sessionId}, \code{correlation\_id},
\code{request\_id}) están presentes en la mayoría de las fuentes de
autenticación, pero al no materializarse como entidades quedaban reducidos a
metadatos de arista, invisibles a la búsqueda estructural. La pregunta
``¿se usó esta misma sesión desde dos direcciones distintas?'' ---evidencia
característica de robo de token o cookie compartida--- no era expresable como
hipótesis y debía responderse con procesamiento ad-hoc sobre los registros
crudos.

El modelo materializa la sesión como una estrella
\[
\mathsf{User} \xrightarrow{\ \mathsf{Uses}\ } \mathsf{Session}
\xrightarrow{\ \mathsf{From}\ } \mathsf{IP}
\]
Una decisión de diseño relevante es que el \emph{rango temporal} de la sesión
(primer evento, último evento, duración, cantidad de eventos) se almacena como
metadato del \textbf{vértice} y no de las aristas: una arista representa un
instante, mientras que una sesión es precisamente aquello que \emph{dura}.
Almacenar un segundo timestamp por arista habría incrementado en un 25\% la
estructura de arista compacta, que es la de mayor volumen del sistema, para
beneficiar únicamente a las pocas relaciones con duración.

Sobre un despliegue real (datos anonimizados; 284.781 eventos de
autenticación), de 60.145 sesiones exactamente una fue utilizada por dos
cuentas distintas, y ambas correspondían a apropiaciones de cuenta confirmadas
de forma independiente. Es la evidencia estructural de operador único más
fuerte del caso, y no depende de correlación temporal alguna.

\subsection{Estructuras y mapeo de logs a grafos}
\textbf{Definición (tripleta de seguridad).} Una tripleta $(s, r, d)$ codifica
entidad fuente $s \in V$, tipo de relación $r \in \Sigma_E$ con marca temporal
$\tau$, y entidad destino $d \in V$. Un parser
$P : \mathrm{Log} \to \mathcal{P}(\mathrm{Tripleta})$ transforma registros de
log en conjuntos de tripletas. El sistema implementa parsers nativos para las
siguientes familias de telemetría:

\begin{itemize}[leftmargin=*]
  \item \textbf{Sysmon:} eventos de Windows Sysmon (creación de procesos,
  conexiones de red, acceso a archivos, consultas DNS, modificaciones de
  registro). Cada tipo de evento (EventID 1, 3, 7, 11, 13, 22) se mapea a
  tripletas específicas.
  \item \textbf{Azure Sentinel:} logs del SIEM cloud de Microsoft, incluyendo
  SecurityEvent, SigninLogs, AzureActivity y CommonSecurityLog. Los EventIDs
  4624/4625 generan tripletas $(\mathsf{User}, \mathsf{Auth}, \mathsf{Host})$.
  \item \textbf{DARPA TC CDM:} formato de proveniencia de los engagements
  Transparent Computing. Requiere resolución en dos pasadas ---catalogar
  entidades por UUID, luego resolver los eventos contra ese catálogo--- porque
  un evento referencia sus extremos por identificadores definidos en otros
  registros. Un evento cuyo extremo no fue catalogado se descarta de forma
  explícita y contabilizada, nunca se completa por inferencia.
  \item \textbf{HAProxy:} logs HTTP de proxy inverso, separando el prefijo
  delimitado por espacios de la línea de petición entrecomillada.
  \item \textbf{FortiGate, Cognito, IIS W3C:} formatos de firewall, identidad
  cloud y servidor web.
  \item \textbf{Genérico JSON y CSV:} registros con campos configurables,
  permitiendo la ingestión de fuentes no estandarizadas.
\end{itemize}

La política de ingesta emplea upsert para vértices (deduplicación por
identificador, fusión de metadatos) e inserción incondicional para aristas,
preservando la granularidad temporal completa. La resolución de entidades
emplea identificadores compuestos ---por ejemplo, \code{user@domain} para
usuarios o \code{host\textbackslash process:pid:timestamp} para procesos---
para desambiguar entidades homónimas como un PID reutilizado en diferentes
sesiones.

\paragraph{Eventos de un solo extremo.} Una relación requiere dos entidades,
de modo que un evento que nombra una sola ---por ejemplo un rechazo de
autenticación que registra la dirección de origen pero ningún usuario--- no
produce arista alguna. En el despliegue mencionado, 642 eventos caían en esta
clase, y entre ellos 399 de los 401 rechazos correspondientes a la ventana de
contención del incidente: precisamente los de mayor valor forense. El sistema
los materializa como self-loops sobre la entidad presente, conservando la
metadata del evento, y reporta la clase por separado en el resumen de ingesta.
La distinción entre ``el parser falló'' y ``el evento nombra un solo extremo''
es material: la primera exige corregir el parser, la segunda es un hecho sobre
los datos.

\subsection{Definiciones sobre hipótesis de caza como consultas de caminos}
El concepto central del sistema es la hipótesis de caza: una secuencia tipada
que describe un patrón de ataque esperado.

\textbf{Hipótesis de caza.} Una secuencia
\[
H = [(v_0, r_1, v_1), (v_1, r_2, v_2), \ldots, (v_{L-1}, r_L, v_L)]
\]
donde $v_i \in \Sigma_V \cup \{*\}$ y $r_j \in \Sigma_E \cup \{*\}$, con $*$
denotando comodín. La hipótesis es consistente si los tipos destino-fuente de
pasos consecutivos coinciden.

\textbf{Camino válido.} Un camino $p = [e_1, e_2, \ldots, e_L]$ sobre aristas
del grafo es válido para $H$ si: (1) los tipos de relación y entidad coinciden
con $H$; (2) se cumple monotonía causal relajada
$\tau(e_{i+1}) \geq \tau(e_i) - \delta$; (3) opcionalmente, una ventana
temporal $[t_{\min}, t_{\max}]$; (4) $k$-simplicidad: ningún vértice aparece
más de $k$ veces; y (5) opcionalmente, hasta $g$ aristas comodín entre pasos
consecutivos.

Esta formalización permite al analista expresar directamente hipótesis
derivadas de MITRE ATT\&CK. Por ejemplo, la hipótesis de movimiento lateral
``un usuario se autentica en un host, desde donde se ejecuta un proceso que se
conecta a un dominio externo'' se codifica como:
\[
H = [(\mathsf{User}, \mathsf{Auth}, \mathsf{Host}),
(\mathsf{Host}, \mathsf{Execute}, \mathsf{Process}),
(\mathsf{Process}, \mathsf{Connect}, \mathsf{Domain})]
\]

\subsection{Tolerancias de despliegue: desfasaje de relojes y huecos}
\label{sec:tolerancias}
La monotonía causal es una restricción natural del dominio: un proceso no puede
conectarse a un servidor C2 antes de ser ejecutado. Resulta además un factor de
poda extremadamente efectivo. Sin embargo, aplicada de forma estricta sobre
telemetría real produce falsos negativos por dos causas distintas, y el modelo
las trata como parámetros declarados en la hipótesis en lugar de supuestos
implícitos.

\paragraph{Desfasaje de relojes ($\delta$).} Una investigación cose eventos
provenientes de relojes que no coinciden: el Sysmon de un endpoint, el registro
de inicio de sesión de un tenant cloud y el flujo CDM de un sensor pueden
diferir en segundos. Una comparación estricta descarta la cadena causal
verdadera cuando el reloj del paso posterior atrasa. El parámetro $\delta$ es
el ancho declarado de ese desacuerdo, no un margen difuso, y se aplica
\textbf{por salto y nunca de forma acumulativa}: cada arista se compara contra
la marca temporal \emph{real} de la anterior, de modo que una cadena de $L$
pasos no puede desviarse más de $\delta$ en ningún salto individual.

\paragraph{Huecos acotados ($g$).} La pérdida de telemetría entre dos acciones
observadas es habitual. El parámetro $g$ admite hasta $g$ aristas comodín
---temporalmente monótonas, honrando $\delta$--- entre dos pasos consecutivos
de $H$. Los vértices puente \textbf{no} forman parte del camino resultante: el
hueco relaja el \emph{matching}, no la \emph{salida}, de modo que el resultado
sigue siendo la cadena de entidades ancla y las operaciones aguas abajo
(deduplicación, agregación, $k$-simplicidad) permanecen inalteradas.

% ======================================================================
\section{Arquitectura del sistema}

\subsection{Visión general}
El sistema se implementa como una aplicación de escritorio con tres capas:
(1) un núcleo en Rust con toda la lógica de dominio (grafo en memoria, parsers,
búsqueda, puntuación endógena de anomalía y bridge GNN para extracción de
características de subgrafos); (2) un backend Tauri v2 que expone el núcleo
como comandos y una API HTTP REST con autenticación Bearer token; y (3) un
frontend React con TypeScript y Cytoscape.js para visualización interactiva. Un
servidor MCP (Model Context Protocol) expone además las operaciones del motor
como herramientas para asistentes de IA.

La elección de un grafo en memoria con estructuras nativas (HashMaps, vectores)
en lugar de una base de datos de grafos externa (Neo4j, ArangoDB) responde al
patrón de acceso del Threat Hunting: se ingieren datos de sesión, se ejecutan
búsquedas interactivas intensivas y se descartan al finalizar, donde el enfoque
en memoria logra tiempos sub-milisegundo frente a los $>100$~ms de una travesía
temporal en Neo4j. El control explícito sobre el algoritmo habilita además
optimizaciones difíciles de expresar en lenguajes declarativos: la monotonía
causal requiere mantener estado sobre la última marca temporal durante la
travesía, lo que en Cypher exigiría procedimientos almacenados. El código
fuente (parsers y motor en Rust) se publica bajo licencia GNU GPLv3 en
\url{https://github.com/Base4Security/GraphHunter}.

\subsection{Estructuras de datos}
El motor mantiene cinco estructuras de datos en memoria:
\begin{enumerate}[leftmargin=*]
  \item \textbf{Entidades} (\code{HashMap<Id, Entity>}): índice principal de
  vértices por identificador único.
  \item \textbf{Adyacencia saliente} en formato CSR (Compressed Sparse Row),
  ordenada por \emph{(fuente, marca temporal)}: lista de adyacencia para
  travesía directa.
  \item \textbf{Índice por tipo de entidad} (\code{HashMap<$\Sigma_V$,
  HashSet<Id>>}): localiza todos los vértices de un tipo dado en $O(1)$,
  fundamental para la fase de semillas del algoritmo.
  \item \textbf{Adyacencia inversa:} para travesía en dirección opuesta y para
  las podas de semi-join bidireccional.
  \item \textbf{Índice secundario de relaciones} (\code{HashMap<(Id,
  $\Sigma_E$), Vec<Relation>>}): recupera en $O(1)$ las aristas salientes de un
  vértice con un tipo de relación específico, evitando el escaneo $O(\deg)$ de
  la lista de adyacencia completa. Impacto especialmente alto en hubs de alto
  grado.
\end{enumerate}

El ordenamiento \emph{(fuente, marca temporal)} de la adyacencia habilita la
optimización estructural más relevante del motor: obtener el rango de sucesores
temporalmente válidos mediante búsqueda binaria, en $O(\log \deg)$ en lugar de
un escaneo $O(\deg)$ por paso.

\subsection{Algoritmo de búsqueda: DFS con poda por dominio}
El algoritmo central, \emph{Temporal-Hypothesis-Search}, opera en dos fases:

\textbf{Fase 1 (semillas):} obtiene del índice de tipos todos los vértices cuyo
tipo coincide con $v_0$ de la hipótesis. Sobre este conjunto se aplican tres
filtros previos: un filtro de etiquetas de vecindario a un salto (NLF), un
filtro de alcanzabilidad etiquetada a $k$ saltos con cotas temporales, y una
poda de semi-join bidireccional que verifica la secuencia completa
tipo/relación sobre bitsets de candidatos por posición. Los tres son
condiciones necesarias: un fallo en la poda es una pérdida de rendimiento,
nunca un error de correctitud.

\textbf{Fase 2 (DFS con backtracking):} para cada semilla, lanza una búsqueda
en profundidad aplicando las siguientes reglas de poda en cada paso:

\begin{enumerate}[leftmargin=*]
  \item \textbf{Poda de tipo de relación:} descarta aristas cuyo tipo no
  coincide con $r_i$ del paso actual. Cuando el índice secundario de relaciones
  está disponible, esta poda se convierte en un pre-filtro $O(1)$.
  \item \textbf{Poda de tipo de entidad destino:} descarta vértices destino
  cuyo tipo no coincide con $v_i$ del paso actual.
  \item \textbf{Poda de monotonía causal:} exige
  $\tau(e_{i+1}) \geq \tau(e_i) - \delta$. Implementada como búsqueda binaria
  sobre la adyacencia ordenada, de modo que la restricción no se evalúa por
  arista sino que define directamente el rango a recorrer.
  \item \textbf{Poda de ventana temporal:} restringe a
  $\tau(e_i) \in [t_{\min}, t_{\max}]$ cuando el analista especifica un período
  de interés. La ventana se aplica en \emph{cada} salto, no solo en el primero.
  \item \textbf{Poda de $k$-simplicidad:} impide que un vértice aparezca más de
  $k$ veces en el camino ($k=1$ por defecto). Permite detectar patrones
  cíclicos (callbacks C2, movimiento lateral con retorno) cuando $k > 1$.
  \item \textbf{Ruptura de simetría:} para cadenas palindrómicas fuerza una
  dirección canónica, reduciendo a la mitad el conjunto de salida sin pérdida
  de información.
\end{enumerate}

En el peor caso teórico (un solo tipo de entidad, un solo tipo de relación, sin
restricción temporal), el problema se reduce a búsqueda de caminos en grafos
dirigidos, que es NP-completo por reducción desde Camino Hamiltoniano Dirigido
(Karp, 1972). Sin embargo, las restricciones del dominio de ciberseguridad
hacen el problema tratable en la práctica:

\textbf{Tratabilidad práctica.} Con $|\Sigma_V| = 10$ y $|\Sigma_E| = 11$, la
poda por tipo de relación retiene una fracción esperada $1/11 \approx 9\%$ de
aristas, y la poda por tipo de entidad retiene otra fracción $1/10$. La
monotonía causal retiene en promedio ${\sim}50\%$ de las aristas restantes. El
factor de ramificación efectivo resulta
\[
b_{\mathrm{eff}} = \bar{d} \cdot \tfrac{1}{11} \cdot \tfrac{1}{10} \cdot
\tfrac{1}{2} \approx \frac{\bar{d}}{220},
\]
que cae por debajo de 1 cuando el grado medio $\bar{d} < 220$ ---condición
satisfecha en todos los grafos de telemetría prácticos. Cuando
$b_{\mathrm{eff}} < 1$, el árbol de búsqueda converge y el tiempo esperado es
$O(|S_0| \cdot c)$ donde $S_0$ son las semillas y $c$ es una constante,
resultando lineal en el tamaño del grafo.

\paragraph{Nota sobre la uniformidad.} Este análisis supone una distribución
uniforme de tipos, supuesto que los datos reales violan: el corpus Mordor
concentra un 85\% de sus nodos en el tipo $\mathsf{Registry}$. La
Sección~\ref{sec:podas} contrasta las fracciones teóricas con las medidas, y la
diferencia es informativa: la poda de tipo de entidad rinde \emph{más} de lo
predicho, no menos.

\subsection{Operadores especializados y su semántica temporal}
\label{sec:operadores}
Junto al DFS, el motor incorpora operadores físicos worst-case-optimal para
formas específicas: un enumerador Leapfrog Triejoin con intersección SIMD para
triángulos y cliques, y un operador Yannakakis para consultas acíclicas. Un
planificador selecciona el operador según la forma de la hipótesis.

La integración de estos operadores expone un riesgo de correctitud que merece
señalarse porque no es evidente y fue observado empíricamente. Una cadena
homogénea de seis pasos $T \xrightarrow{r} T$ es la forma habitual de expresar
``un camino de seis saltos''; es además el número de aristas dirigidas de un
4-clique. Un planificador que despache según la \emph{forma} sola enruta esa
hipótesis al enumerador de cliques, que responde una pregunta distinta de la
formulada. El sistema exige por ello un marcador explícito del autor además de
la forma: la forma no es la intención.

De modo análogo, la semántica temporal del enumerador debe coincidir con la del
DFS en cuatro aspectos que resultaron divergir en una revisión sistemática:
(i) la comparación de monotonía debe ser no estricta, dado que las fuentes
marcan eventos causalmente ordenados dentro del mismo segundo; (ii) el borde
superior de la ventana debe ser inclusivo; (iii) la ventana debe aplicarse a
todas las aristas del patrón y no solo a la primera ---restringirla a la
primera \emph{sobre-reporta}, emitiendo coincidencias que la ventana del
analista excluía---; y (iv) ante aristas paralelas, el operador debe preguntar
si \emph{alguna} combinación respeta el orden temporal, y no evaluar únicamente
la marca mínima de cada par. Este último punto es el de mayor impacto práctico:
las aristas paralelas entre el mismo par de entidades son la norma en
telemetría real, y colapsarlas al mínimo descarta coincidencias legítimas de
forma rutinaria. La verificación se realiza contra una implementación de
referencia independiente, precisamente porque comparar variantes que comparten
la misma premisa solo detecta errores de codificación, no premisas erróneas.

\subsection{Puntuación y priorización de amenazas}
Para guiar al analista en la formulación de hipótesis y la priorización de
resultados, el sistema calcula tres métricas de centralidad:

\textbf{Centralidad de grado normalizada:}
\[
\mathrm{score}_{\deg}(v) = \frac{\deg_{\mathrm{in}}(v) +
\deg_{\mathrm{out}}(v)}{2|V|}.
\]
Identifica entidades altamente conectadas (hubs potencialmente comprometidos).

\textbf{PageRank temporal:} incorpora decaimiento exponencial por antigüedad de
arista $w(e) = \exp(-\lambda(\tau_{\max} - \tau(e)))$. Prioriza entidades
involucradas en actividad reciente. Implementación: iteración de potencia con
damping $\alpha = 0{,}85$ (Page et al., 1999).

\textbf{Betweenness centrality:} algoritmo de Brandes (2001), exacto en
$O(|V| \cdot |E|)$ o aproximado con $k$ fuentes muestreadas. Identifica
entidades que actúan como intermediarios en múltiples caminos ---nodos de
pivote potenciales en cadenas de movimiento lateral.

El puntaje compuesto $S = w_1 \cdot \deg + w_2 \cdot \mathrm{PR} + w_3 \cdot
\mathrm{BC}$ (con pesos por defecto 0,3, 0,4 y 0,3) permite al analista
identificar rápidamente las entidades más relevantes para iniciar una
investigación, complementando el enfoque hipótesis-dirigido con señales
automáticas de anomalía.

\subsection{Puntuación endógena de anomalía por camino}
Complementariamente a la centralidad por nodo, el sistema implementa una
puntuación de anomalía que opera sobre caminos completos retornados por la
búsqueda. El objetivo es priorizar automáticamente los caminos que involucran
entidades y relaciones estadísticamente inusuales, sin inteligencia de amenazas
externas. El puntaje compuesto $S(p) \in [0,1]$ combina cinco componentes con
pesos configurables: (1) \emph{Entity Rarity} (ER): entidades observadas pocas
veces reciben puntaje alto, calculado como
$\mathrm{ER}(v) = 1 - \ln(\mathrm{freq}(v)) / \ln(f_{\max})$; (2) \emph{Edge
Rarity} (EdgeR): conexiones únicas entre pares de entidades; (3)
\emph{Neighborhood Concentration} (NC): basado en la entropía de Shannon
normalizada del vecindario ---un nodo que solo se conecta con entidades de un
solo tipo sugiere infraestructura como servidores C2; (4) \emph{Temporal
Novelty} (TN): entidades que aparecieron recientemente en el dataset; (5)
\emph{GNN Threat}: clasificación de subgrafos $k$-hop mediante redes neuronales
de grafos.

Una precisión metodológica relevante: la rareza de entidad se normaliza contra
la entidad más frecuente \emph{del mismo tipo}, y solo cuando el tipo cuenta
con una población suficiente de pares; por debajo de ese umbral se recurre al
denominador global. Sin esta corrección, la entidad más activa de un tipo poco
poblado recibe un puntaje de rareza artificialmente alto.

Un modelo ONNX clasifica el vecindario de cada entidad. La retroalimentación es
bidireccional: los puntajes endógenos (ER, TN, NC) alimentan el vector de
características de la GNN como dimensiones de entrada, y la salida de la GNN
alimenta el puntaje compuesto. Esto crea un ciclo cerrado donde la detección
estadística y la clasificación aprendida se potencian mutuamente. La puntuación
también guía la poda durante la búsqueda DFS: nodos con puntaje de anomalía
bajo pueden ser descartados tempranamente mediante una poda admisible de tipo
A\textsuperscript{*}, que solo descarta ramas que \emph{probadamente} no entran
en el top-$k$.

\paragraph{Guarda de precisión.} Un patrón que coincide con una fracción
elevada de la población no discrimina, y un conjunto de resultados grande puede
leerse erróneamente como un hallazgo importante. El sistema mide la prevalencia
sobre \emph{entidades semilla} ---no sobre caminos, cuyo número se infla con el
fan-out--- y marca el patrón como no discriminante por encima de un umbral
configurable. La guarda se abstiene explícitamente cuando el conjunto de
resultados fue truncado por el límite de materialización: un conjunto acotado
es un piso, no una medición.

\subsection{Integración con agentes de IA vía MCP}
El sistema expone una API HTTP REST en localhost con autenticación Bearer
token, permitiendo que herramientas externas consulten el grafo sin usar la
interfaz gráfica. Sobre esta API, un servidor MCP (Model Context Protocol)
traduce las operaciones del motor en herramientas invocables por asistentes de
IA. Las primitivas expuestas están orientadas al flujo de trabajo de
investigación: búsqueda de entidades, expansión de vecindarios $k$-hop,
ejecución de hipótesis con ranking por anomalía, obtención de detalles de nodos
con puntajes GNN, y creación de notas de análisis. Es requisito de diseño que
toda operación disponible en la interfaz gráfica lo esté también por API y por
MCP, de modo que un analista pueda delegar tareas exploratorias a un agente que
opera directamente sobre el grafo de la sesión activa, manteniendo al humano en
el loop para decisiones de alto nivel.

% ======================================================================
\section{Evaluación}

\subsection{Entorno experimental}
Todas las mediciones utilizan compilaciones en modo release de Rust sobre un
solo core (AMD Ryzen 5 PRO 5650U, 16~GB RAM). No se emplea paralelismo en la
fase de búsqueda para obtener mediciones reproducibles.

\textbf{Datasets reales.} Tres datasets de escala creciente:
\begin{itemize}[leftmargin=*]
  \item \textbf{APT Simulation} (30 eventos Sysmon): cadena APT completa desde
  acceso inicial (phishing) hasta exfiltración, con procesos encadenados y
  conexiones de red. Simula técnicas T1566, T1059, T1071, T1041 de MITRE
  ATT\&CK.
  \item \textbf{Sentinel Simulation} (32 eventos Azure Sentinel): escenario de
  compromiso cloud-to-on-premises con autenticaciones Azure AD, actividad en
  recursos cloud y movimiento lateral hacia infraestructura local.
  \item \textbf{Mordor Combined} (963.000 eventos): telemetría real del
  proyecto Mordor (Rodriguez, 2020), con sesgo extremo de tipos (85\% nodos
  Registry) y distribución de grado power-law con $d_{\max} = 9.411$. Este
  dataset estresa particularmente las podas de tipo y la gestión de hubs de
  alto grado.
\end{itemize}

\textbf{Reproducibilidad.} Las mediciones de escala se obtienen mediante un
arnés incorporado al binario del motor, invocable como un único comando, con
semilla fija y salida estructurada (CSV o JSON). El arnés distingue
explícitamente lo que mide ---latencia, escala y una estimación estructural de
memoria--- de lo que no mide: la eficacia de detección requiere etiquetas de
verdad de campo y se aborda por separado.

\subsection{Caso de estudio 1: detección de movimiento lateral}
Consideremos un escenario de movimiento lateral mediante Pass-the-Hash, una
técnica frecuente en APTs donde un atacante compromete un endpoint vía
phishing, extrae credenciales con herramientas como Mimikatz, y las utiliza
para autenticarse en servidores internos donde ejecuta comandos de
exfiltración.

\textbf{Flujo en SIEM tradicional.} Tras detectar un pico de autenticaciones
fallidas seguido de una exitosa desde una IP inusual, reconstruir la cadena
causal exige una secuencia de consultas independientes ---IP origen en
firewall, usuario y procesos asociados, conexiones de red del proceso, hosts y
autenticaciones posteriores--- cada una con extracción manual de identificadores
del resultado anterior.

\textbf{Flujo con el enfoque propuesto.} La investigación se codifica como una
sola hipótesis:
\begin{align*}
H = [\,&(\mathsf{IP}, \mathsf{Connect}, \mathsf{Host}),
(\mathsf{Host}, \mathsf{Auth}, \mathsf{User}),\\
&(\mathsf{User}, \mathsf{Execute}, \mathsf{Process}),
(\mathsf{Process}, \mathsf{Write}, \mathsf{File})\,]
\end{align*}
El algoritmo retorna todos los caminos válidos con monotonía causal en
milisegundos. El analista visualiza la cadena completa: IP atacante $\to$ Host
comprometido $\to$ Usuario admin $\to$ \code{Mimikatz.exe} $\to$ archivo
exfiltrado.

\subsection{Caso de estudio 2: ransomware fileless}
En un escenario de ransomware fileless, un documento de Office ejecuta una
macro que invoca PowerShell (T1059.001). PowerShell descarga un payload en
memoria (T1105), inyecta código en \code{svchost.exe} (T1055), realiza
consultas DNS a un dominio C2 (T1071.004), y cifra archivos (T1486). La
hipótesis correspondiente:
\[
H = [(\mathsf{Process}, \mathsf{Spawn}, \mathsf{Process}),
(\mathsf{Process}, \mathsf{DNS}, \mathsf{Domain}),
(\mathsf{Process}, \mathsf{Write}, \mathsf{File})]
\]
La ventaja del grafo se manifiesta especialmente en la detección de inyección
de código. En logs lineales, la relación entre PowerShell y \code{svchost.exe}
se pierde entre miles de eventos de spawn de procesos.

\subsection{Caso de estudio 3: compromiso cloud-to-on-premises}
Un escenario emergente en entornos híbridos: un atacante obtiene credenciales
de Azure AD mediante phishing, accede a recursos cloud, y desde allí pivota
hacia infraestructura on-premises. La hipótesis:
\[
H = [(\mathsf{User}, \mathsf{Auth}, \mathsf{Host}),
(\mathsf{Host}, \mathsf{Connect}, \mathsf{IP}),
(\mathsf{IP}, \mathsf{Auth}, \mathsf{Host})]
\]
El parser de Sentinel normaliza tanto los SigninLogs (Azure AD) como los
SecurityEvent (Windows on-premises) al mismo modelo de grafo, permitiendo
trazar la cadena completa a través de ambos entornos en una sola consulta.

\subsection{Caso de estudio 4: continuidad de sesión}
La materialización de la sesión (Sección~\ref{sec:sesion}) habilita una clase
de hipótesis inexpresable en el modelo original:
\[
H = [(\mathsf{Session}, \mathsf{From}, \mathsf{IP})] \quad \text{con}
\quad |\{\mathsf{IP}\}| \geq 3
\]
El umbral es 3 y no 2 de forma deliberada: los proxies inversos y las capas
\emph{backend-for-frontend} aparecen como segunda dirección en prácticamente
toda sesión, de modo que dos direcciones es el caso normal y no el hallazgo.
Sobre el despliegue anonimizado, de 60.145 sesiones 12.695 involucraban más de
una dirección; de ellas 12.652 exactamente dos ---ruido de infraestructura--- y
solo 43 (el 0,07\% del total) tres o más. Elevar el umbral de 2 a 3 convierte
un conjunto de resultados inmanejable en uno revisable a mano, y ese conjunto
contenía los casos de apropiación confirmados. El umbral no se derivó de un
principio sino de medir la prevalencia del patrón antes de atribuirle
significado.

\subsection{Resultados de rendimiento}
El dataset Mordor, con 963K eventos originales comprimidos a 772 nodos y 21.963
aristas, presenta el caso más exigente por su distribución power-law
($d_{\max} = 9.411$). Aun así, la búsqueda con $L = 2$ completa en 4,3~ms
visitando solo 13 nodos frente a 22.735 del baseline sin poda.

\begin{table}[htbp]
\centering
\caption{Rendimiento de la búsqueda sobre datasets reales.}
\begin{tabular}{lrrrrrr}
\toprule
\textbf{Dataset} & $|V|$ & $|E|$ & $L$ & \textbf{Tiempo} & \textbf{Speedup} &
\textbf{Resultados}\\
\midrule
APT      & 24  & 30     & 2 & 0,005 ms & $7{,}9\times$    & 0\\
APT      & 24  & 30     & 5 & 0,010 ms & $4{,}6\times$    & 7\\
Sentinel & 28  & 32     & 3 & 0,006 ms & $2{,}1\times$    & 3\\
Mordor   & 772 & 21.963 & 2 & 4,3 ms   & $1{,}747\times$  & 438\\
Mordor   & 772 & 21.963 & 3 & 4,4 ms   & $5{,}121\times$  & 42\\
Mordor   & 772 & 21.963 & 5 & 4,4 ms   & $11{,}600\times$ & 0\\
\bottomrule
\end{tabular}
\end{table}

\subsection{Escalabilidad: barrido a densidad constante}
\label{sec:escala}
Las mediciones anteriores caracterizan datasets concretos, pero no aíslan el
efecto del tamaño. Para ello se ejecuta un barrido sobre grafos
Erd\H{o}s--Rényi generados con \emph{grado medio fijo} en 10, de modo que
$p = \bar{d}/(n-1)$ mantiene la densidad constante y cada fila aísla el efecto
de $|E|$ en lugar de confundirlo con una densificación progresiva. Se ejecutan
15 repeticiones por punto con semilla fija, hipótesis de $L = 3$.

\begin{table}[htbp]
\centering
\caption{Latencia de búsqueda frente a $|E|$, a densidad constante
($\bar{d} = 10$, $L = 3$).}
\begin{tabular}{rrrr}
\toprule
$|V|$ & $|E|$ & \textbf{p50} & \textbf{p95}\\
\midrule
1.000    & 9.925     & 133 $\mu$s & 215 $\mu$s\\
5.000    & 49.522    & 841 $\mu$s & 1,3 ms\\
10.000   & 99.870    & 1,9 ms     & 12 ms\\
25.000   & 250.613   & 7,5 ms     & 9,8 ms\\
50.000   & 500.185   & 18 ms      & 22 ms\\
100.000  & 999.037   & 33 ms      & 48 ms\\
250.000  & 2.498.604 & 121 ms     & 134 ms\\
\bottomrule
\end{tabular}
\end{table}

Un incremento de $251\times$ en aristas (de 9.925 a 2.498.604) produce un
incremento de $908\times$ en la mediana de latencia, lo que corresponde a un
crecimiento aproximado de $O(|E|^{1{,}2})$: ligeramente superlineal, y no
atribuible a densificación puesto que el grado medio permanece fijo. El cuello
de botella es el trabajo de búsqueda y no la memoria ---la representación CSR
consume del orden de 32 bytes por arista, de modo que $2{,}5$ millones de
aristas ocupan menos de 100~MB--- ni la emisión de resultados, cuyo costo
resultó despreciable frente al recorrido.

Extrapolando esta tendencia al objetivo de un engagement DARPA TC E3
($\sim$5,9 millones de aristas), una consulta de tres saltos se ubicaría en el
orden de los 340~ms: en el límite de lo interactivo para una consulta
individual, y del orden de segundos para un catálogo completo de 20--30
hipótesis.

\subsection{Efectividad de las reglas de poda}
\label{sec:podas}
\begin{table}[htbp]
\centering
\caption{Porcentaje de candidatos eliminados por cada regla de poda.}
\begin{tabular}{lrr}
\toprule
\textbf{Regla de poda} & \textbf{Teórico} & \textbf{Mordor}\\
\midrule
Tipo de relación (pre-filtro por índice) & 90,9\% & ${\sim}0\%$\textsuperscript{*}\\
Tipo de entidad destino                  & 90,0\% & 92,3\%\\
Monotonía causal                         & 50,0\% & 38,1\%\\
$k$-simplicidad                          & $<1\%$ & 0\%\\
\bottomrule
\end{tabular}

\vspace{2mm}
{\footnotesize \textsuperscript{*}Pre-filtrada por índice secundario antes del
bucle de podas.}
\end{table}

La poda de tipo de relación muestra 0\% en el bucle porque el índice secundario
la resuelve antes de la iteración ---las aristas con tipo incorrecto nunca
entran al bucle de podas. La poda de tipo de entidad es el filtro dominante en
ejecución, eliminando $>90\%$ de candidatos en datos reales. La monotonía
causal aporta un 38\% adicional en Mordor, menor que el 50\% teórico, lo que
refleja la distribución temporal no uniforme de eventos reales de seguridad.

\subsection{Análisis comparativo: SIEM vs. grafo}
La diferencia más significativa es el escalamiento: en SIEM, el tiempo de
investigación crece linealmente con la longitud de la cadena de ataque,
mientras que en el enfoque propuesto el tiempo de ejecución es esencialmente
constante respecto a $L$ (dado que $b_{\mathrm{eff}} < 1$). Para
investigaciones complejas con $L \geq 4$, la reducción supera un orden de
magnitud.

\begin{table}[htbp]
\centering
\caption{Comparativa operativa para investigación de movimiento lateral
($L = 4$).}
\begin{tabular}{lll}
\toprule
\textbf{Métrica} & \textbf{SIEM tradicional} & \textbf{Enfoque propuesto}\\
\midrule
Consultas/operaciones     & 4 secuenciales             & 1 hipótesis\\
Tiempo de formulación     & ---                        & 1--2 min\\
Tiempo de ejecución       & 8--20 min total            & $<5$ ms\\
Extracción manual de IDs  & Sí, en cada paso           & No\\
Descubrimiento exhaustivo & No (3--5 caminos)          & Sí (todos)\\
Escalamiento con $L$      & Lineal (${\sim}L\times$2--5 min) & Constante\\
Visualización causal      & No nativa                  & Grafo interactivo\\
\bottomrule
\end{tabular}
\end{table}

% ======================================================================
\section{Discusión}

\subsection{Resolución de las barreras operativas}
El sistema resuelve las tres barreras de la Sección 1.2. \textbf{(i) Pivoteo
iterativo.} Las $L$ consultas manuales secuenciales de un SIEM se reemplazan
por una única hipótesis tipada que el DFS recorre automáticamente respetando
tipo y temporalidad, conectando IPs, PIDs y usuarios de forma estructural; el
MTTI cae de 8--20 minutos a 1--2 minutos para cadenas de 4 pasos.
\textbf{(ii) Sobrecarga cognitiva.} El modelo mental que el analista debía
sostener en memoria de trabajo se externaliza: el grafo almacena las relaciones
causales como aristas tipadas, la visualización las hace navegables, y los
puntajes de anomalía priorizan los nodos de interés, mitigando el cuello de
botella de los $7 \pm 2$ ítems de Miller (1956). \textbf{(iii) Fragmentación
del contexto.} Los parsers multi-formato ingestan telemetría heterogénea bajo
una ontología única, la deduplicación de identidad reconstruye el contexto de
cada entidad a través de fuentes, y la búsqueda por hipótesis retorna solo los
subgrafos que coinciden con el patrón.

\textbf{Descubrimiento exhaustivo.} La enumeración revela rutas que el pivoteo
manual no encontraría: en Mordor la hipótesis con $L = 2$ retorna 438 caminos
válidos, frente a los 3--5 que un analista explora típicamente.

\textbf{Ciclo cerrado detección-investigación.} La clasificación GNN como
componente de la puntuación de anomalía unifica la detección automática y la
investigación dirigida por hipótesis sobre el mismo grafo: ante un subgrafo
anómalo, el analista formula una hipótesis para investigar el contexto causal
sin exportar datos ni reconstruir el grafo en otra herramienta.

\subsection{Limitaciones}
\textbf{Escalabilidad en memoria.} El tamaño del dataset queda limitado por la
RAM disponible; para volúmenes mayores se recurre a compactación temporal
(${\sim}40\%$ de reducción en aristas) o ventanas deslizantes. La medición de
la Sección~\ref{sec:escala} indica que el límite práctico actual es de trabajo
de búsqueda antes que de memoria.

\textbf{Visualización a escala.} Por encima de ${\sim}500$ nodos la
visualización directa pierde utilidad (``hairball''); se mitiga con la búsqueda
por hipótesis, BFS acotada y los puntajes de anomalía.

\textbf{Dependencia de la hipótesis.} El método es tan bueno como las hipótesis
formuladas; un analista sin experiencia en ATT\&CK podría no dar con las
correctas, lo que se complementa con los puntajes de anomalía y motiva la
generación automática de hipótesis como trabajo futuro.

\textbf{Semántica de los operadores especializados.} Como se detalla en la
Sección~\ref{sec:operadores}, los operadores worst-case-optimal exigen una
alineación explícita de su semántica temporal con la del DFS. Al momento de
escribir, dicha alineación está verificada para los enumeradores de triángulos
y de 4-cliques contra implementaciones de referencia independientes; los
operadores de 5-cliques y 6-ciclos permanecen desactivados por defecto hasta
completar la misma verificación. Se documenta esta limitación de forma
explícita porque un operador más rápido que responde una pregunta distinta es
peor que no tenerlo.

\textbf{Betweenness sobre proyección estática.} La betweenness opera sobre la
proyección estática del multigrafo, ignorando el orden temporal y pudiendo
sobreestimar nodos cuyas rutas no son temporalmente factibles; computar
betweenness temporal de camino más temprano es \#P-hard (Buss et al., 2022),
por lo que la aproximación estática es un compromiso deliberado.

\textbf{Dependencia de modelos GNN pre-entrenados.} El componente GNN Threat
requiere un modelo ONNX entrenado sobre grafos de proveniencia etiquetados, y
su calidad depende de la representatividad del dataset respecto al SOC. Se
mitiga con el peso configurable $w_5 = 0$ (desactivado por defecto) y el ajuste
fino sobre datos propios. Los modelos entrenados sobre corpus de proveniencia
deben además publicar sus salvedades: un clasificador cuya verdad de campo
etiqueta archivos afectados por un ataque ---y no los procesos que lo
ejecutaron--- es en la práctica un detector binario, con independencia del
número de clases de su cabeza de salida.

% ======================================================================
\section{Conclusiones y trabajo futuro}
Este trabajo demuestra que la modelización de telemetría de seguridad como
grafo temporal tipado, combinada con búsqueda de hipótesis por DFS con reglas
de poda, reduce significativamente el tiempo y la complejidad cognitiva de las
investigaciones de Threat Hunting. Las restricciones inherentes al dominio de
ciberseguridad ---tipado de entidades y relaciones, causalidad temporal---
convierten un problema teóricamente intratable en uno con tiempo esperado
lineal en la práctica.

Los resultados principales son: (1) aceleraciones de $1{,}747\times$ a
$11{,}600\times$ sobre baseline sin poda en datasets reales; (2) un
crecimiento de latencia cercano a $O(|E|^{1{,}2})$ medido a densidad constante
hasta $2{,}5$ millones de aristas, con mediana de 121~ms; (3) reducción de un
orden de magnitud en el MTTI (de 8--20 minutos a 1--2 minutos para
investigaciones de 4 pasos); (4) descubrimiento exhaustivo de caminos de ataque
versus la exploración parcial del pivoteo manual; (5) integración de
clasificación GNN sobre subgrafos $k$-hop como componente de la puntuación de
anomalía; y (6) la incorporación de tolerancias de despliegue ---desfasaje de
relojes y huecos acotados--- que convierten supuestos implícitos sobre la
calidad de la telemetría en parámetros declarados de la consulta.

Una observación metodológica transversal merece registrarse: varias de las
limitaciones corregidas durante este trabajo ---eventos de un solo extremo
invisibles al grafo, sesiones no materializadas, divergencia semántica entre
operadores--- no fueron anticipadas en el diseño sino descubiertas usando la
herramienta contra incidentes reales. Esto sugiere que el catálogo de supuestos
de un sistema de este tipo se construye más confiablemente en despliegue que en
el pizarrón.

\textbf{Líneas futuras.} Se identifican: búsqueda en streaming sobre ventanas
deslizantes con eventos fuera de orden (Lai et al., 2022); generación
automática de hipótesis a partir de los puntajes GNN y embeddings del grafo,
complementando el catálogo ATT\&CK; integración con feeds de inteligencia de
amenazas (IoCs vía STIX/TAXII); evaluación de usabilidad con analistas de SOC
mediante métricas como NASA-TLX; ordenamiento de variables sensible a
cardinalidad para las formas no lineales de hipótesis; completar la alineación
semántica de los operadores worst-case-optimal restantes; y una arquitectura
distribuida que separe el fast-path en endpoint (inferencia NPU) del servidor
de investigación (grafo temporal + DFS) para entornos enterprise.

% ======================================================================
\section*{Referencias}
\small
\begin{list}{}{\setlength{\leftmargin}{1.5em}\setlength{\itemindent}{-1.5em}
\setlength{\itemsep}{2pt}}

\item Aberger, C. R., Lamb, A., Tu, S., Nötzli, A., Olukotun, K. y Ré, C.
(2017). EmptyHeaded: A relational engine for graph processing. \emph{ACM
Transactions on Database Systems}, 42(4), 1--44.

\item Brandes, U. (2001). A faster algorithm for betweenness centrality.
\emph{Journal of Mathematical Sociology}, 25(2), 163--177.

\item Buss, S., Molter, H., Niedermeier, R. y Renken, M. (2022). Computing
optimal temporal walks under waiting-time constraints. \emph{Algorithmica}, 84,
2754--2802.

\item Cai, T., Li, C. y Pei, J. (2023). Exact subgraph matching via
edge-ordered temporal graphs. En \emph{Proceedings of IEEE ICDE} (pp.
733--744). IEEE.

\item Han, X., Pasquier, T., Bates, A., Mickens, J. y Seltzer, M. (2020).
Unicorn: Runtime provenance-based detector for advanced persistent threats. En
\emph{Proceedings of NDSS}. Internet Society.

\item Hossain, M. N., Milajerdi, S. M., Wang, J., Eshete, B., Gjomemo, R.,
Sekar, R., Stoller, S. y Venkatakrishnan, V. N. (2017). SLEUTH: Real-time
attack scenario reconstruction from COTS audit data. En \emph{Proceedings of
USENIX Security Symposium} (pp. 487--504). USENIX.

\item Hutchins, E. M., Cloppert, M. J. y Amin, R. M. (2011).
Intelligence-driven computer network defense informed by analysis of adversary
campaigns and intrusion kill chains. En \emph{Proceedings of the 6th
International Conference on Information Warfare and Security}. Academic
Publishing International.

\item Jia, Z., Xiong, Y., Nan, Y., Zhang, Y., Zhao, J. y Wen, M. (2024). MAGIC:
Detecting advanced persistent threats via masked graph representation learning.
En \emph{Proceedings of the 33rd USENIX Security Symposium} (pp. 5197--5214).
USENIX Association.

\item Karp, R. M. (1972). Reducibility among combinatorial problems. En R. E.
Miller y J. W. Thatcher (Eds.), \emph{Complexity of Computer Computations} (pp.
85--103). Plenum Press.

\item Kim, H., Lee, Y. y Shin, K. (2022). Temporal subgraph isomorphism. En
\emph{Proceedings of IEEE ICDE} (pp. 532--543). IEEE.

\item Kipf, T. N. y Welling, M. (2017). Semi-supervised classification with
graph convolutional networks. En \emph{Proceedings of ICLR}. OpenReview.

\item Lai, S., Li, X., Lu, J., Luo, C. y Zhao, D. (2022). Scalable temporal
subgraph pattern matching. \emph{Proceedings of the VLDB Endowment}, 15(11),
2647--2660.

\item Manzoor, E., Milajerdi, S. M. y Akoglu, L. (2016). Fast memory-efficient
anomaly detection in streaming heterogeneous graphs. En \emph{Proceedings of
the 22nd ACM SIGKDD International Conference on Knowledge Discovery and Data
Mining} (pp. 1035--1044). ACM.

\item Milajerdi, S. M., Gjomemo, R., Eshete, B., Sekar, R. y Venkatakrishnan,
V. N. (2019a). HOLMES: Real-time APT detection through correlation of
suspicious information flows. En \emph{Proceedings of IEEE Symposium on
Security and Privacy} (pp. 1137--1152). IEEE.

\item Milajerdi, S. M., Eshete, B., Gjomemo, R. y Venkatakrishnan, V. N.
(2019b). Poirot: Aligning attack behavior with kernel audit records for cyber
threat hunting. En \emph{Proceedings of ACM CCS} (pp. 1795--1812). ACM.

\item Miller, G. A. (1956). The magical number seven, plus or minus two: Some
limits on our capacity for processing information. \emph{Psychological Review},
63(2), 81--97.

\item Page, L., Brin, S., Motwani, R. y Winograd, T. (1999). \emph{The PageRank
citation ranking: Bringing order to the web} (Technical Report 1999-66).
Stanford InfoLab.

\item Robbins, A., Schroeder, W. y Foss, L. (2017). An ACL-based attack path
analysis for Active Directory. BloodHound: Six Degrees of Domain Admin.
SpecterOps.

\item Rodriguez, R. (2020). \emph{Mordor: Pre-recorded security events from
simulated adversarial techniques}. Open Threat Research.
\url{https://github.com/OTRF/mordor}

\item Veldhuizen, T. L. (2014). Leapfrog Triejoin: A simple, worst-case optimal
join algorithm. En \emph{Proceedings of the 17th International Conference on
Database Theory} (pp. 96--106).

\item Wang, Q., Hassan, W. U., Li, D., Jee, K. y Yu, X. (2020). You are what
you do: Hunting stealthy malware via data provenance analysis. En
\emph{Proceedings of NDSS}. Internet Society.

\item Wang, Q., Hassan, W. U., Li, D., Jee, K. y Yu, X. (2022). ThreaTrace:
Detecting and tracing host-based threats in node level through provenance graph
learning. \emph{IEEE Transactions on Dependable and Secure Computing}, 19(6),
3792--3808.

\item Yannakakis, M. (1981). Algorithms for acyclic database schemes. En
\emph{Proceedings of the 7th International Conference on Very Large Data
Bases} (pp. 82--94).

\end{list}
\normalsize

\end{document}
